\section{La concordancia como predicción selectiva}
\label{sec:disc_concordancia}
% Extensión estimada: 2 páginas, cinco párrafos (ficha 6.3).
% Fuentes: sec:marco_pred_selectiva, sec:marco_concordancia, sec:met_ruteo; sec:res_gen_concordancia (1 204 afirmados,
% 102 errores afirmados, 230 errores de Swin-Tiny, 248 diferenciales, 51 derivados), sec:res_gen_rag (171 y 72),
% tab:res_dermlip (kNN 0.781), sec:res_gen_pipeline, sec:res_melanoma, sec:res_calibracion, sec:res_gen_estado_arte.
% Cifras derivadas: cobertura 1 204 / 1 503 = 0.80; riesgo 102 / 1 204 = 0.085; error global 230 / 1 503 = 0.153;
% ambos equivocados 230 - 72 = 158; esperable si fueran independientes 0.153 x 0.219 x 1 503 = 50.
% Reemplaza la parte del ruteo de las antiguas 6.3 y 6.4.

El ruteo por concordancia es, en los términos de la Sección~\ref{sec:marco_pred_selectiva}, un clasificador selectivo:
Swin-Tiny propone el diagnóstico, y la concordancia con el \gls{knn} decide si se afirma. Leído así, sus resultados
(Sección~\ref{sec:res_gen_concordancia}) se resumen en dos cifras. El \emph{pipeline} afirma el diagnóstico en el
80\,\% de las imágenes de prueba, que es su cobertura, y entre ellas el clasificador se equivoca en el 8.5\,\%, que es
su riesgo selectivo, frente al 15.3\,\% de error sobre la prueba completa. Es decir, renunciar a afirmar una de cada
cinco imágenes reduce a algo más de la mitad la proporción de diagnósticos afirmados erróneos. Además, a diferencia de
un umbral sobre la probabilidad, la regla no tiene ningún parámetro ajustado sobre la prueba, de modo que esa
reducción no está sobreestimada por la selección.

Que la concordancia funcione depende de que los dos predictores se equivoquen en imágenes distintas. La
Sección~\ref{sec:marco_concordancia} lo planteaba como una hipótesis: un clasificador que ajusta fronteras entre las
clases y un \gls{knn} que vota con los casos más parecidos, sobre una representación aprendida de otra manera, deberían
tener sesgos distintos. Los resultados la respaldan solo en parte. Los predictores difieren en muchas imágenes: el
\gls{knn} acierta 72 en que Swin-Tiny se equivoca, y Swin-Tiny acierta 171 en que el \gls{knn} se equivoca
(Sección~\ref{sec:res_gen_rag}). Esos desacuerdos son los que el ruteo aprovecha, y por eso en el nivel diferencial el
par de candidatos contiene la clase verdadera en el 84.7\,\% de los casos.

Sus errores, sin embargo, coinciden mucho más de lo que coincidirían por azar. Los dos predictores se equivocan a la
vez en 158 imágenes, y en 102 de ellas eligen la misma clase errónea. Si sus errores fueran independientes, con una
tasa de error de 0.153 para Swin-Tiny y de 0.219 para el \gls{knn}, se equivocarían a la vez en unas 50 imágenes, de
modo que la coincidencia observada es unas tres veces mayor. Una explicación posible es que los dos predictores
aprenden de las mismas 6\,447 imágenes de entrenamiento: el clasificador ajusta sus fronteras con ellas y el \gls{knn}
las usa como casos de referencia. Si una lesión de prueba se parece a lesiones de entrenamiento de otra clase, o si
esas lesiones están mal etiquetadas, los dos heredan la misma confusión, por distintas que sean sus representaciones.
Otra explicación, compatible con la anterior, es que una parte de esas imágenes sea ambigua para cualquier modelo,
como lo es la distinción entre melanoma y nevo para los especialistas. En cualquier caso, la heterogeneidad de las
representaciones no basta para que los errores sean independientes. Otras señales de confianza, como la probabilidad de
Swin-Tiny, que está bien calibrada cuando se entrena sin pesos por clase (Sección~\ref{sec:res_calibracion}), no
intervienen en el ruteo, y no se evaluó si detectan una parte de esos errores.

Los errores compartidos se concentran, además, donde el costo clínico es mayor. El \emph{pipeline} se eligió y se
evaluó por la consistencia diagnóstica, que con la cascada coincide con la exactitud del clasificador y está dominada,
por lo tanto, por los nevos, que son dos tercios de la prueba. En el melanoma, la prioridad clínica es la
sensibilidad, y ahí el \emph{pipeline} queda más débil: casi uno de cada cuatro melanomas se afirma como otra clase, la
mayoría como nevo (Sección~\ref{sec:res_gen_concordancia}). El ruteo mitiga parte del problema, porque cuando un
melanoma pasa al nivel diferencial el par de candidatos casi siempre lo incluye, pero no puede mitigar los melanomas
que los dos predictores confunden. Al mismo tiempo, el umbral fijado en validación, que lleva la detección de
Swin-Tiny de tres a cuatro de cada cinco melanomas (Sección~\ref{sec:res_melanoma}), no se usa en el ancla. Es decir,
el \emph{pipeline} optimiza una métrica que trata igual todos los errores, mientras que la clínica no los trata igual,
y la configuración que maximiza la consistencia no tiene por qué ser la preferible en la práctica.

Lo que sí logra el ruteo es convertir parte de los errores en una duda explícita. En el nivel diferencial, el reporte
presenta dos candidatos en lugar de afirmar el diagnóstico equivocado del clasificador, y en la derivación no se emite
reporte: se evitan 20 reportes incorrectos y dos contradictorios, a costa de 29 diagnósticos correctos. El valor de ese
intercambio depende de cuánto cueste en la práctica un error frente a una consulta con el especialista, una relación
que este trabajo no midió. El único trabajo revisado que combina un ensamble con la derivación,
el de \textcite{khan2025ensemble}, no informa cuántos casos deriva ni con qué acierto (Sección~\ref{sec:res_gen_estado_arte}),
de modo que la cobertura y el riesgo selectivo de este trabajo no tienen todavía con qué compararse. En conjunto, el
\emph{pipeline} puede ordenar la atención del especialista, señalando qué casos revisar primero, pero no reemplazarla,
la misma conclusión a la que llegan \textcite{jiang2026multimodal} sobre los \glspl{vlm} que evalúan en consultas
reales.
